\section{控制混杂偏倚}

每当我们着手评估一个因素（$X$）对另一个因素（$Y$）的效应时，都会出现这样的问题：是否应当针对某些其他因素（$Z$）的可能变异，对我们的测量作调整（或“标准化”）——这些因素又称“协变量”“伴随变量”或“混杂因子”（Cox 1958, p.~48）。调整相当于把总体划分为相对于 $Z$ 同质的若干组，在每一同质组内评估 $X$ 对 $Y$ 的效应，再对结果取平均（如 \textup{(3.13)} 那样）。这种调整的虚幻性质早在 1899 年就已被认识：当时 Karl Pearson 发现了如今所谓的辛普森悖论（见 \S 6.1）——两个变量之间的任何统计关系，都可能因在分析中纳入额外因素而被逆转。例如，我们可能发现吸烟的学生比不吸烟的学生成绩更好；但若对年龄作调整，则在每个年龄组中吸烟者成绩更低；再对家庭收入作调整，则在每个收入—年龄组中吸烟者又比不吸烟者成绩更好；如此等等。

尽管经过一个世纪的分析，辛普森的逆转仍在“困住不谨慎者”（Dawid 1979），而它所提出的实际问题——对给定协变量的调整是否恰当——一直未能得到数学处理。例如，流行病学家至今仍在争论“混杂”的含义（Grayson 1987; Shapiro 1997），并常常对错误的协变量集合作调整（Weinberg 1993；另见第六章）。Rosenbaum and Rubin (1983) 与 Pratt and Schlaifer (1988) 的潜在结局分析引出了一个名为“可忽略性”的概念，它用反事实词汇重述了混杂问题，却未能为研究者提供一条可操作的准则来指导协变量的选择（见 \S 11.3.2）。可忽略性表述为：“若给定 $Z$ 时，‘$Y$ 在 $X$ 为 $x$ 时会取的值’与 $X$ 独立，则 $Z$ 是一个可容许的协变量集。”由于反事实不可观测，而关于反事实之条件独立的判断又难以从普通科学知识中断言，问题便一直悬而未决：究竟应当用什么样的准则来判定哪些变量适于调整？

\S 3.3.1 用因果图这一友好的语言，给出调整问题的一个一般而形式的解。\S 3.3.2 把这一结果推广到受 $X$ 影响的非标准协变量，这类协变量需要多步调整。最后，\S 3.3.3 用一个例子说明这些准则的用法。

\subsection{后门准则}

假定给定因果图 $G$，以及 $G$ 中可观测变量子集 $V$ 上的非实验数据；并假定我们希望估计干预 $\doxp(X=x)$ 会对一组响应变量 $Y$ 产生什么效应，其中 $X$ 与 $Y$ 是 $V$ 的两个子集。换言之，我们希望在给定 $G$ 所编码之假设的条件下，由 $P(v)$ 的样本估计来估计 $P(y\mid\hat{x})$。

我们将说明存在一个简单的图形检验，Pearl (1993b) 称之为“后门准则”，它可直接施于因果图，以检验变量集 $Z\subseteq V$ 是否足以识别 $P(y\mid\hat{x})$。\footnote{这一准则也可由 Spirtes 等 (1993) 的定理 7.1 得到。另一种只用一个 d-分离检验的准则，将在 \S 3.4 中建立（见 \textup{(3.37)}）。}

\begin{definition}[后门]
若变量集 $Z$ 关于 DAG $G$ 中的有序变量对 $(X_i,X_j)$ 满足：
\begin{enumerate}[label=(\roman*),leftmargin=2.2em]
\item $Z$ 中没有任何节点是 $X_i$ 的后代；且
\item $Z$ 阻断了 $X_i$ 与 $X_j$ 之间所有含有指向 $X_i$ 之箭头的路径，
\end{enumerate}
则称 $Z$ 满足\emph{后门准则}。类似地，若 $X$ 与 $Y$ 是 $G$ 中两个不相交的节点子集，则当 $Z$ 关于任一满足 $X_i\in X$、$X_j\in Y$ 的数对 $(X_i,X_j)$ 都满足该准则时，称 $Z$ 关于 $(X,Y)$ 满足后门准则。
\end{definition}

“后门”之名呼应条件 (ii)：它要求只阻断那些箭头指向 $X_i$ 的路径；这些路径可视为经由“后门”进入 $X_i$。例如在图 3.4 中，集合 $Z_1=\{X_3,X_4\}$ 与 $Z_2=\{X_4,X_5\}$ 满足后门准则，而 $Z_3=\{X_4\}$ 不满足，因为 $X_4$ 未能阻断路径 $(X_i,X_3,X_1,X_4,X_2,X_5,X_j)$。第十一章（\S 11.3.1）给出 (i) 与 (ii) 的直观说明。

\begin{theorem}[后门调整]
若变量集 $Z$ 关于 $(X,Y)$ 满足后门准则，则 $X$ 对 $Y$ 的因果效应可识别，且由下式给出：
\begin{equation}
P(y\mid\hat{x}) = \sum_{z} P(y\mid x,z)\,P(z). \tag{3.19}
\end{equation}
\end{theorem}

\textup{(3.19)} 中的求和代表对 $Z$ 作调整所得到的标准公式；使 \textup{(3.19)} 中等式成立的变量 $X$，在 Rosenbaum and Rubin (1983) 中被称为“给定 $Z$ 时条件可忽略的”。把可忽略性条件化归为定义 3.3.1 的图形准则，就把关于反事实依赖的判断替换为关于因果关系的普通判断，正如图中所示。该图形准则可由系统程序检验，适用于任何大小与形状的图。该准则还使分析者能够搜索最优的协变量集——即使用户测量成本或抽样变异性最小的集合 $Z$（Tian et al. 1998）。把类似的图形准则用于识别线性结构方程中的路径系数，将在第五章演示。在流行病学研究中的应用见 Greenland 等 (1999a) 与 Glymour and Greenland (2008)，那里集合 $Z$ 被称为“充分集”；“可容许集”或“去混杂集”是更好的术语。

\begin{figure}[htbp]
\centering
\begin{tikzpicture}[causal,scale=0.95,every node/.style={transform shape}]
  \node[vtx] (X1) at (0.6,2.6) {$X_1$};
  \node[vtx] (X2) at (3.6,2.4) {$X_2$};
  \node[vtx] (X3) at (-0.9,1.7) {$X_3$};
  \node[vtx] (X4) at (1.4,1.5) {$X_4$};
  \node[vtx] (X5) at (4.6,1.6) {$X_5$};
  \node[vtx] (Xi) at (-0.9,0.2) {$X_i$};
  \node[vtx] (X6) at (1.6,0.2) {$X_6$};
  \node[vtx] (Xj) at (4.0,0.2) {$X_j$};
  \draw[ed] (X1)--(X3); \draw[ed] (X1)--(X4);
  \draw[ed] (X4)--(X3); \draw[ed] (X4)--(Xi);
  \draw[ed] (X2)--(X4); \draw[ed] (X2)--(X5);
  \draw[ed] (X5)--(Xj);
  \draw[ed] (X6)--(Xi); \draw[ed] (X6)--(Xj);
  \draw[ed] (Xi)--(Xj);
\end{tikzpicture}
\caption{表示后门准则的图；对变量 $\{X_3,X_4\}$（或 $\{X_4,X_5\}$）作调整可得到 $P(x_j\mid\hat{x}_i)$ 的一致估计。对 $\{X_4\}$ 或 $\{X_6\}$ 作调整则得到有偏估计。}
\label{fig:3.4}
\end{figure}

\noindent\textbf{定理 3.3.2 的证明}\quad Pearl (1993b) 最初给出的证明基于这样一个观察：当 $Z$ 阻断了从 $X$ 到 $Y$ 的一切后门路径时，设定 $X=x$ 与对 $X=x$ 取条件对 $Y$ 的效果相同。这一点最好从图 3.2 的增广图 $G'$ 看出——那里添加了干预弧 $F_X\to X$。若从 $X$ 到 $Y$ 的一切后门路径都被阻断，则从 $F_X$ 到 $Y$ 的一切路径都必须经过 $X$ 的子节点，而若我们对 $X$ 取条件，这些路径便会被阻断。其含义是：给定 $X$ 时 $Y$ 与 $F_X$ 独立，
\begin{equation}
P(y\mid x, F_X=\doxp(x)) = P(y\mid x, F_X=\text{idle}) = P(y\mid x), \tag{3.20}
\end{equation}
也就是说，观测 $X=x$ 与干预 $F_X=\doxp(x)$ 无法区分。

形式上，我们可以按 \textup{(3.9)} 用增广概率函数 $P'$ 写出 $P(y\mid\hat{x})$，并对 $Z$ 取条件，得到
\begin{equation}
P(y\mid\hat{x}) = P'(y\mid F_x) = \sum_{z}P'(y\mid z,F_x)P'(z\mid F_x) = \sum_{z}P'(y\mid z,x,F_x)P'(z\mid F_x). \tag{3.21}
\end{equation}
最后一式中加上 $x$ 是由蕴含式 $F_x\Rightarrow X=x$ 所许可的。为从 \textup{(3.21)} 右端的两个项中消去 $F_x$，我们援引定义 3.3.1 的两个条件。既然 $F_x$ 由若干根节点构成、其子节点仅限于 $X$，它必定与 $X$ 的一切非后代独立，包括 $Z$。于是条件 (i) 给出
\[
P'(z\mid F_x)=P'(z)=P(z).
\]
再援引后门条件 (ii) 与 \textup{(3.20)}，便可在 \textup{(3.21)} 中消去 $F_x$，从而证明 \textup{(3.19)}。\hfill$\blacksquare$

\S 11.3.3 提供了另一种证明。

\subsection{前门准则}

定义 3.3.1 的条件 (i) 反映了“伴随观测不应受处理影响”这一流行做法（Cox 1958, p.~48）。本节说明，受处理影响的伴随变量也可以用来促进因果推断。由此浮现的准则在 Pearl (1995a) 中称为\emph{前门准则}，它将成为识别因果效应之一般检验的第二个构件（\S 3.4 与定理 3.6.1）。

考虑图 3.5 的图：它表示当变量 $X_1,\dots,X_5$ 未被观测、且 $\{X_i,X_6,X_j\}$ 分别重标记为 $\{X,Z,Y\}$ 时的图 3.4 模型。尽管 $Z$ 不满足任何后门条件，对 $Z$ 的测量仍可使 $P(y\mid\hat{x})$ 得到一致估计。我们将通过把 $P(y\mid\hat{x})$ 的表达式化简为可由被观测分布函数 $P(x,y,z)$ 计算的公式来表明这一点。

\begin{figure}[htbp]
\centering
\begin{tikzpicture}[causal,scale=1.0,every node/.style={transform shape}]
  \node[vtx] (X) at (0,1.8) {$X$};
  \node[vtx] (Z) at (0,0.6) {$Z$};
  \node[vtx] (Y) at (0,-0.7) {$Y$};
  \node[vtx] (U) at (1.5,1.8) {$U$};
  \draw[ed] (U)--(X); \draw[ed] (U) to[bend left=25] (Y);
  \draw[ed] (X)--(Z); \draw[ed] (Z)--(Y);
\end{tikzpicture}
\caption{表示前门准则的图。对 $Z$ 作两步调整可得到 $P(y\mid\hat{x})$ 的一致估计。}
\label{fig:3.5}
\end{figure}

与图 3.5 相联系的联合分布可分解（式 \textup{(3.5)}）为
\begin{equation}
P(x,y,z,u) = P(u)P(x\mid u)P(z\mid x)P(y\mid z,u). \tag{3.22}
\end{equation}
由 \textup{(3.10)}，干预 $\doxp(x)$ 移除因子 $P(x\mid u)$，并诱导出干预后分布
\begin{equation}
P(y,z,u\mid\hat{x}) = P(y\mid z,u)P(z\mid x)P(u). \tag{3.23}
\end{equation}
对 $z$ 与 $u$ 求和便得到
\begin{equation}
P(y\mid\hat{x}) = \sum_{z}P(z\mid x)\sum_{u}P(y\mid z,u)P(u). \tag{3.24}
\end{equation}
为从 \textup{(3.24)} 右端消去 $u$，我们利用图 3.5 所编码的两个条件独立假设：
\begin{equation}
P(u\mid z,x) = P(u\mid x), \tag{3.25}
\end{equation}
\begin{equation}
P(y\mid x,z,u) = P(y\mid z,u). \tag{3.26}
\end{equation}
由此得到等式
\begin{equation}
\sum_{u}P(y\mid z,u)P(u) = \sum_{x}\sum_{u}P(y\mid z,u)P(u\mid x)P(x) = \sum_{x}P(y\mid x,z)P(x), \tag{3.27}
\end{equation}
并使 \textup{(3.24)} 化简为只涉及被观测量的形式：
\begin{equation}
P(y\mid\hat{x}) = \sum_{z}P(z\mid x)\sum_{x'}P(y\mid x',z)P(x'). \tag{3.28}
\end{equation}
\textup{(3.28)} 右端的一切因子都可由非实验数据一致地估计，因而 $P(y\mid\hat{x})$ 亦可估计。于是，只要我们找到一个满足 \textup{(3.25)} 与 \textup{(3.26)} 条件的中介变量 $Z$，便掌握了 $X$ 对 $Y$ 之因果效应的无偏非参数估计量。

式 \textup{(3.28)} 可以解释为后门公式的两步应用。第一步，我们求 $X$ 对 $Z$ 的因果效应；由于图 3.5 中没有从 $X$ 到 $Z$ 的未阻断后门路径，我们简单地有
\[
P(z\mid\hat{x})=P(z\mid x).
\]
第二步，我们计算 $Z$ 对 $Y$ 的因果效应；由于存在从 $Z$ 到 $Y$ 的后门路径 $Z\leftarrow X\leftarrow U\to Y$，它不再能等同于条件概率 $P(y\mid z)$。然而，既然 $X$ 阻断了（d-分离了）这条路径，$X$ 就可以在后门准则中充当伴随变量，从而使我们能够按 \textup{(3.19)} 计算 $Z$ 对 $Y$ 的因果效应，得到 $P(y\mid\hat{z})=\sum_{x'}P(y\mid x',z)P(x')$。最后，我们通过
\[
P(y\mid\hat{x})=\sum_z P(y\mid\hat{z})P(z\mid\hat{x})
\]
把两个因果效应结合起来，它化简为 \textup{(3.28)}。

在形式地定义假设之后，我们用定理来总结这一结果。

\begin{definition}[前门]
若变量集 $Z$ 关于有序变量对 $(X,Y)$ 满足：
\begin{enumerate}[label=(\roman*),leftmargin=2.2em]
\item $Z$ 截断了从 $X$ 到 $Y$ 的一切有向路径；
\item 没有从 $X$ 到 $Z$ 的未阻断后门路径；且
\item 从 $Z$ 到 $Y$ 的一切后门路径都被 $X$ 阻断，
\end{enumerate}
则称 $Z$ 满足\emph{前门准则}。
\end{definition}

\begin{theorem}[前门调整]
若 $Z$ 关于 $(X,Y)$ 满足前门准则，且 $P(x,z)>0$，则 $X$ 对 $Y$ 的因果效应可识别，并由下式给出：
\begin{equation}
P(y\mid\hat{x}) = \sum_{z}P(z\mid x)\sum_{x'}P(y\mid x',z)P(x'). \tag{3.29}
\end{equation}
\end{theorem}

定义 3.3.3 所述的条件过于严格；条件 (ii)、(iii) 所排除的某些后门路径其实可以允许，只要它们被某些伴随变量阻断即可。例如，图 3.1 中的变量 $Z_2$ 就由于 $Z_1$ 阻断了从 $X$ 到 $Z_2$ 以及从 $Z_2$ 到 $Z_3$ 的一切后门路径，而关于 $(X,Z_3)$ 满足一个类前门的准则。为分析这类错综结构、包括后门条件与前门条件的嵌套组合，\S 3.4 将引入一套更强的符号机械，以绕开诸如导出 \textup{(3.28)} 时所用的代数操作。但首先让我们看一个例子，说明前门条件的可能应用。

\subsection{例：吸烟与基因型理论}

考虑关于吸烟（$X$）与肺癌（$Y$）之间关系的那场百年之争（Spirtes et al. 1993, pp.~291--302）。据许多人说，烟草业正是靠如下论点阻止了反吸烟立法：吸烟与肺癌之间的观测相关，可以由某种致癌基因型（$U$）来解释，而这种基因型包含对尼古丁的先天渴求。

肺部焦油沉积量（$Z$）是一个有望满足定义 3.3.3 所列条件的变量，因而契合图 3.5 的结构。为满足条件 (i)，我们必须假定吸烟除经由焦油沉积之外，对肺癌的产生没有影响。为满足条件 (ii) 与 (iii)，我们必须假定：即便某基因型加剧了肺癌的产生，它除间接地（经由吸烟）之外对肺部焦油量也没有影响。同样，我们必须假定没有其他影响焦油沉积的因素会对吸烟有任何影响。最后，定理 3.3.4 的条件 $P(x,z)>0$ 要求：肺部高焦油水平不仅是吸烟的结果，也可以是其他因素（如暴露于环境污染物）的结果；并且某些吸烟者可能没有焦油（也许由于极其高效的排焦油机制）。最后一个条件是否满足，可以在数据中检验。

为演示我们如何评估吸烟在多大程度上增加（或减少）肺癌风险，我们假定有一项假想研究，其中 $X,Y,Z$ 三个变量在一个大样本、随机抽取的人群中被同时测量。为简化阐述，我们进一步假定三个变量都是二值的，取真（1）或假（0）。表 3.1 给出了一项关于焦油、癌症与吸烟之关系研究的假想数据集。它显示，95\% 的吸烟者与 5\% 的不吸烟者肺部出现了高焦油水平。此外，有焦油沉积的受试者中 81\% 患肺癌，而无焦油沉积者中仅 14\%。最后，在这两组（有焦油与无焦油）的每一组内，吸烟者的癌症比例都远高于不吸烟者。

\begin{table}[htbp]
\centering
\caption{假想研究数据：$P(x,z)$ 与 $P(Y=1\mid x,z)$}
\label{tab:3.1}
\small
\begin{tabular}{llcc}
\toprule
组别 & 组类型 & 组规模（占总体 \%） & 组内癌症病例 \% \\
\midrule
$X=0,\ Z=0$ & 不吸烟者，无焦油 & 47.5 & 10 \\
$X=1,\ Z=0$ & 吸烟者，无焦油   & 2.5  & 90 \\
$X=0,\ Z=1$ & 不吸烟者，有焦油 & 2.5  & 5  \\
$X=1,\ Z=1$ & 吸烟者，有焦油   & 47.5 & 85 \\
\bottomrule
\end{tabular}
\end{table}

这些结果似乎证明吸烟是肺癌的主要成因。然而烟草业可以辩称，该表讲的是另一个故事——吸烟实际上降低了患肺癌的风险。其论证如下。若你决定吸烟，则你形成焦油沉积的机会是 95\%，而不吸烟时只有 5\%。为评估焦油沉积的效应，我们分别考察吸烟者与不吸烟者两组。表中显示，焦油沉积在两组中都有保护作用：在吸烟者中，焦油沉积使癌症率从 90\% 降到 85\%；在不吸烟者中，使癌症率从 10\% 降到 5\%。因此，无论我是否天生渴求尼古丁，我都应当去追求肺部焦油沉积的保护作用，而吸烟提供了获得这些沉积的极有效手段。

为解决两种解释之间的争端，我们现在把前门公式（式 \textup{(3.29)}）用于表 3.1 的数据。我们希望计算一位随机抽取者在以下两种行动之一下患癌症的概率：吸烟（设 $X=1$）或不吸烟（设 $X=0$）。

代入 $P(z\mid x)$、$P(y\mid x,z)$ 与 $P(x)$ 的相应取值，我们有
\begin{equation}
\begin{aligned}
P(Y=1\mid\doxp(X=1)) &= .05(.10\times.50+.90\times.50)+.95(.05\times.50+.85\times.50)\\
&= .05\times.50+.95\times.45=.4525,
\end{aligned} \tag{3.30}
\end{equation}
\begin{equation}
\begin{aligned}
P(Y=1\mid\doxp(X=0)) &= .95(.10\times.50+.90\times.50)+.05(.05\times.50+.85\times.50)\\
&= .95\times.50+.05\times.45=.4975.
\end{aligned}
\end{equation}
于是，与预期相反，数据证明吸烟对健康略有裨益。

表 3.1 的数据显然不合实际，是刻意构造出来支持基因型理论的。然而这项练习的目的，是演示关于机制运作的合理定性假设、与结合非实验数据，如何能够产生对因果效应的精确量化评估。在现实中，我们预期涉及中介变量的观测研究会否证基因型理论——例如表明吸烟的中介后果（如焦油沉积）在吸烟者与不吸烟者中都会增加而非降低癌症风险。那时便可用 \textup{(3.29)} 的估计量来量化吸烟对癌症的因果效应。
